\documentclass[11pt,a4paper]{article}
\usepackage{amsmath}

\begin{document}

\section{Equations}
\label{sec:eq}
Energy is conserved, so $E = mc^2$ holds for a particle at rest, and
a fraction such as $\frac{a}{b}$ sits inside the line. Equation~\eqref{eq:euler}
is Euler's identity:
\begin{equation}
  e^{i\pi} + 1 = 0
  \label{eq:euler}
\end{equation}
and the next display is not numbered:
\[
  \sum_{k=1}^{n} k = \frac{n(n+1)}{2}.
\]
Aligned equations line up at the relation:
\begin{align}
  f(x) &= x^2 + 2x + 1 \label{eq:f} \\
       &= (x+1)^2. \label{eq:g}
\end{align}
Section~\ref{sec:eq} has equations \eqref{eq:f} and~\eqref{eq:g}, and
\begin{equation*}
  \alpha + \beta \leq \sqrt{\gamma^2 + \delta^2}
\end{equation*}
ends it.

\end{document}
